%% ma: L10-B13-0081
%% lop: 10
%% chuong: 5
%% bai: 13
%% dang: 10.13.10
%% loai: TN4
%% muc: TH
%% dap_an: "C"
%% nguon: "chuyen-de-cac-so-dac-trung-cua-mau-so-lieu-khong-ghep-nhom-toan-10-chuong-trinh-moi.pdf (GV. Phan Nhật Linh) (P03), Bài 02 Dạng 2 Phần I Câu 6"
%% kien_thuc: "Lớp 10 Bài 13 (trung vị, tứ phân vị)"
%% trang_thai: chinh-thuc
%% ghi_chu: "Tự duyệt theo yêu cầu giáo viên 22/10/2026."
\begin{ex}
Số xe ô tô một cửa hàng bán được mỗi tháng trong năm 2021 được ghi lại ở bảng sau:
\begin{center}
\begin{tabular}{|c|c|c|c|c|c|c|c|c|c|c|c|c|}\hline
Tháng & $1$ & $2$ & $3$ & $4$ & $5$ & $6$ & $7$ & $8$ & $9$ & $10$ & $11$ & $12$\\\hline
Số xe & $45$ & $28$ & $31$ & $34$ & $32$ & $35$ & $37$ & $33$ & $33$ & $35$ & $34$ & $37$\\\hline
\end{tabular}
\end{center}
Tứ phân vị thứ nhất $Q_1$ của mẫu số liệu là
\choice{$31{,}5$}{$32$}{\True $32{,}5$}{$34$}
\loigiai{Sắp xếp: $28$; $31$; $32$; $33$; $33$; $34$; $34$; $35$; $35$; $37$; $37$; $45$.

Nửa dưới: $28$; $31$; $32$; $33$; $33$; $34$ nên $Q_1=\dfrac{32+33}2=32{,}5$. Chọn C.}
\end{ex}
